\documentclass[10pt]{letter}

\usepackage[T1]{fontenc}
\usepackage[utf8]{inputenc}
\usepackage[margin=0.82in]{geometry}
\usepackage{hyperref}
\usepackage{enumitem}

\hypersetup{hidelinks}
\setlength{\parskip}{0.45em}
\setlength{\parindent}{0pt}

\begin{document}
\pagestyle{empty}

\begin{flushright}
Yuhao Peng\\
School of Mathematics and Physics\\
China University of Geosciences (Wuhan)\\
Wuhan, Hubei 430074, China\\
\href{mailto:pngyuo@cug.edu.cn}{pngyuo@cug.edu.cn}
\end{flushright}

\vspace{0.5em}
9 November 2026

\vspace{0.5em}
Dear Editors,

\textbf{Re: Submission of ``Stage-conditional shift normalization for
subject-robust EEG sleep staging''}

We submit the enclosed manuscript for consideration as a research article in
\emph{Neurocomputing}.

Sleep stage classifiers lose accuracy when they are transferred to a new
patient or clinical center. The common label-free response is to replace the
batch-normalization statistics of the target recording. That operation
implicitly treats subject shift as one recording-wide offset. Our diagnostic
results show that this view is incomplete: after marginal mean correction, a
subject probe refitted inside each sleep stage still performs well above
chance, whereas the ordinary overall probe is dominated by the removable
offset. The finding locates a stage-conditional component of subject
dependence that global normalization cannot address.

We introduce stage-conditional shift normalization (SCSN). SCSN estimates a
separate first-moment shift for each sleep stage from the model's own soft
posterior and subtracts the posterior-weighted shift before classification.
It is closed form, requires no target labels, optimizer, or weight update, and
contains the conventional marginal correction as one endpoint of a
one-parameter family. The method applies unchanged to batch- and
group-normalized stagers and leaves the backbone and classifier frozen.

The study evaluates SCSN on Sleep-EDF and the independent HMC Sleep Staging
Database, using two backbones and two normalization architectures. The key
results are:

\begin{enumerate}[leftmargin=1.2em,itemsep=0.08em,topsep=0.2em]
\item SCSN reduces the within-stage subject probe from 0.74 to 0.47 on
Sleep-EDF EEGNet and from 0.53 to 0.21 on HMC EEGNet;
\item accuracy is preserved on Sleep-EDF and improves by 1.4--1.9 points on
HMC;
\item on group-normalized models, SCSN improves accuracy significantly in
three of four settings and worst-patient conformal coverage in all four;
\item matched-statistic, mismatched-statistic, hard-label, low-rank, and
classifier-null-space controls isolate the mechanism;
\item cross-cohort transfer identifies the practical boundary of the method:
a weak target posterior reduces SCSN to marginal correction.
\end{enumerate}

The manuscript reports these qualifications rather than presenting SCSN as
universally superior. Code, evaluation scripts, and released result summaries
will be available at \url{https://github.com/YuhaoPeng1103/SleepSCSN}.

The manuscript is original, is not under consideration elsewhere, and has not
been published previously. Both authors approved the submission. The authors
declare no competing interests, received no specific funding, suggest no
reviewers, and request the subscription publication route.

Thank you for considering the manuscript.

\vspace{0.8em}
Sincerely,

\vspace{0.8em}
Yuhao Peng\\
On behalf of Zhiwen Song and Yuhao Peng

\end{document}
